\section{Removing the Lie fields from twenty-nine through forty-three}
\label{sec:lie-small-fields}

\begin{theorem}[D-BI-20: the remaining Lie-field bound]
\label{D-BI-20}
Consider a finite-cover Weil sixfold source of exact imaginary
quadratic character field, with nonsplit actual polarization and a
period map of rank nine.  Let $\tau$ be an absolutely irreducible
inducing block of degree $e\ge10$, and suppose that its finite image
$J$ is primitive and tensor indecomposable in the factorwise sense
specified in D-BI-19.  Suppose the simple factors of the normal core
in D-BI-19 are nonalternating groups of Lie type over $\mathbf F_q$.
Then
\[
                              q\le27.                       \tag{BI.72}
\]
The inducing character field may be larger than the original
imaginary quadratic field.  Tensor-factor permutations and all
outer automorphisms compatible with the representation are allowed.
The conclusion concerns a hypothetical dominant source; it does not
construct sources for the remaining small fields.
\end{theorem}

\begin{proof}
By D-BI-19 there is a normal irreducibly acting central product
$E=E_1\cdots E_t$, whose quasisimple factors are permuted
transitively by $J$.  Write $S=E_i/Z(E_i)$.  Restriction to the
commuting factors expresses the representation of $E$ as a tensor
product of irreducible representations of the $E_i$.  Their degrees
are equal, say $a$, because conjugation by $J$ permutes the factors
transitively.  Thus
\[
                              e=a^t.                        \tag{BI.73}
\]
Each factor representation is faithful: a factor element acting
trivially there acts trivially on the full tensor product, and $E_i$
is an actual subgroup of $\operatorname{GL}(V)$.  In particular its
degree is at least the minimum degree of a nontrivial projective
complex representation of $S$.

The case $q\ge47$ is already excluded by D-BI-18 and D-BI-19,
because $e\ge10$.  The prime powers between $28$ and $46$ are
precisely
\[
                    \mathcal Q=\{29,31,32,37,41,43\}.
\]
For these six values, the exact integer inequality in D-BI-18 gives
the following bounds on the \emph{total} degree:
\[
\begin{array}{c|rrrrrr}
q&29&31&32&37&41&43\\ \hline
e\le U(q)&87&45&33&21&15&15.
\end{array}                                                   \tag{BI.74}
\]

We first remove all types other than $A_1$.  The published
cross-characteristic projective lower bounds for the classical
groups, as collected in \cite[Table~1, author version p.~14]{BI:Hoffman},
give, for every classical simple group over one of these fields
other than $\operatorname{PSL}_2(q)$,
\[
                        a\ge(q^2-1)/2\ge420.                \tag{BI.75}
\]
Here ``other than $A_1$'' includes the rank-two symplectic and
unitary cases.  To check the coarse bound directly from that table:
the linear minimum is at least $q^2+q-1$; the odd-characteristic
symplectic minimum is $(q^n-1)/2$, $n\ge2$; its
even-characteristic counterpart is
$(q^n-1)(q^n-q)/(2(q+1))$, $n\ge2$; the unitary minima are
$(q^{2n}-1)/(q+1)$, $n\ge2$, and
$(q^{2n+1}-q)/(q+1)$, $n\ge1$.  Each is at least the right
side of (BI.75).  The three orthogonal rows, with their stated
minimum ranks, are still larger.  All exceptional parameter pairs
in those classical rows have field size at most four, so none
occurs here.  We use only the published classical rows of this
source, not its rows attributed to an unpublished exceptional-group
manuscript.

For exceptional types we instead use the characteristic-zero
minimum-degree theorem and complete tables of
\cite[Theorem~2.1 and \S5]{BI:Luebeck}.  There are no exceptional
Schur multipliers for these six field sizes: the complete list in
\cite[Chapter~5, Table~1, numbered p.~35]{BI:HHM} has none.
Consequently the simply connected character tables used by
L\"ubeck cover every possible quasisimple central cover here.
For the types $G_2,{}^3D_4,F_4,{}^2F_4,E_6,{}^2E_6,E_7,E_8$,
the first nontrivial degree in the respective tables is at least
$q(q-1)>87$ at the parameters under consideration.  This coarse
inequality can also be checked using the first-degree expressions
retained in \texttt{check\_lie\_small\_fields.py}.
The only possible Suzuki parameter in $\mathcal Q$ is $q=32$,
where the minimum is
\[
                         (32-1)\sqrt{32/2}=124.
\]
The small Ree groups require $q$ to be an odd power of three, and
so do not occur in $\mathcal Q$.  Thus every exceptional type
also contradicts (BI.74).  These are minimum-degree inputs from
published classification tables, not results of a new enumeration
of all finite groups.

We are left with $S=\operatorname{PSL}_2(q)$.  Its projective
minimum-degree bound, from the first row of the same classical
table, is
\[
 a\ge a_0(q):=\frac{q-1}{\gcd(2,q-1)}.
\]
The relevant numbers are
\[
\begin{array}{c|rrrrrr}
q&29&31&32&37&41&43\\ \hline
a_0(q)&14&15&31&18&20&21\\
a_0(q)^2&196&225&961&324&400&441.
\end{array}                                                   \tag{BI.76}
\]
Every entry in the second row exceeds its corresponding $U(q)$.
Equations (BI.73)--(BI.76) therefore force $t=1$.  In particular
$E$ itself acts irreducibly, rather than merely one tensor factor
of a higher-degree representation.  The values $q=41,43$ are
already impossible because $a_0(q)>U(q)$.

For precision about the remaining covers, the usual universal
central cover of $\operatorname{PSL}_2(q)$ is
$\operatorname{SL}_2(q)$, except for the exceptional multipliers
at $q=4,9$.  The exceptional list and the distinction between
ordinary Lie-type covers and exceptional covers are explicitly
recorded in \cite[Introduction, numbered p.~4; Chapter~5,
Table~1, numbered p.~35]{BI:HHM}.  Hence at the parameters here
the group $E$ is $\operatorname{SL}_2(q)$ or its quotient
$\operatorname{PSL}_2(q)$; for even $q$ these coincide.  We must
therefore check inverse conjugacy and automorphisms on
$\operatorname{SL}_2(q)$, not only on its projective quotient.

At $q=37$, all irreducible degrees of $\operatorname{SL}_2(37)$
belong to
\[
                       \{1,18,19,36,37,38\}.               \tag{BI.77}
\]
This is the specialization $\alpha=q-1$, $\beta=1$, $f=1$ of
\cite[Theorem~B, p.~4]{BI:WhiteSL2}; degrees of the projective
quotient are a subset.  Directly imposing
$12\lceil e(37-9)/(4\cdot37)\rceil\le2e+6$ with
$3\le e\le21$ leaves only
\[
                         \{3,4,5,9,10,15,21\}.
\]
This set is disjoint from (BI.77), so $q=37$ is excluded by the
exact ceiling inequality, without a self-duality shortcut.

We next prove that every element of $\operatorname{SL}_2(q)$
is conjugate to its inverse when $q$ is even or $q\equiv1\pmod4$.
A noncentral semisimple matrix $x$ and $x^{-1}$ have the same
characteristic polynomial and are conjugate in
$\operatorname{GL}_2(q)$.  The determinant map on the
$\operatorname{GL}_2(q)$-centralizer of $x$ is surjective onto
$\mathbf F_q^*$: for a split torus this is immediate from the
diagonal matrices, and for a nonsplit torus it is the surjective
norm $\mathbf F_{q^2}^*\to\mathbf F_q^*$.  Multiplying a
conjugator by a suitable centralizer element therefore makes its
determinant one.  Thus the semisimple inverse conjugacy takes
place in $\operatorname{SL}_2(q)$.

For odd $q$, every remaining noncentral element is conjugate to
$\pm(I+bE_{12})$ for some $b\ne0$.  If $q\equiv1\pmod4$,
choose $u\in\mathbf F_q$ with $u^2=-1$.  Conjugation by
$\operatorname{diag}(u,u^{-1})$ replaces $b$ by $-b$, producing
the inverse.  For even $q$, a nontrivial unipotent is an
involution.  Central elements are already their own inverses.
This proves the assertion on the full special linear group, and
it descends to its central quotients.  Its class description agrees
with \cite[\S3, pp.~5--7]{BI:WhiteSL2}.

It follows that every irreducible character of $E$ is real when
$q=29$ or $32$.  Indeed $\chi(x^{-1})=\overline{\chi(x)}$
for a finite-group character, and inverse-conjugate elements have
the same character value.  A real irreducible character is
self-dual and has a nondegenerate invariant complex bilinear form,
possibly alternating.  D-BI-18a applies to the entire normalizer
$J$, since the $E$-invariant form is unique up to scalar.  Its
hypotheses hold here: $q\ge19$, $e\ge10$, the inducing block is
primitive, and the original actual polarization is nonsplit.
Therefore $q=29,32$ are excluded.

Finally let $q=31$.  The outer automorphism group of
$\operatorname{PSL}_2(31)$ has order two, generated by the
diagonal automorphism; there is no field automorphism because
$31$ is prime.  This is the specialization of
\cite[\S2, p.~3; \S3, p.~6]{BI:WhiteSL2}.  The same assertion
holds for $\operatorname{SL}_2(31)$.  To justify the latter
point, an automorphism inducing the identity on the central
quotient has the form $x\mapsto z(x)x$, with $z$ a homomorphism
to the center.  The group $\operatorname{SL}_2(31)$ is perfect,
so $z=1$.  Perfectness here can be seen directly: commutators of
$\operatorname{diag}(u,u^{-1})$ with upper and lower unipotents,
for $u^2\ne1$, produce all those unipotents, and those matrices
generate $\operatorname{SL}_2(31)$.  Thus its automorphisms
inject into those of its projective quotient, and the diagonal
automorphism supplies the nontrivial outer class.

This diagonal automorphism is class-inverting on
$\operatorname{SL}_2(31)$.  For semisimple elements it preserves
the characteristic polynomial, hence the special-linear conjugacy
class by the centralizer argument above; that class also contains
the inverse.  It fixes the central elements.  On each of the two
unipotent square classes, conjugation by
$\operatorname{diag}(v,1)$ for a nonsquare $v$ exchanges the
square classes.  Inversion has the same effect because $-1$ is a
nonsquare in $\mathbf F_{31}$.  This also treats the negatives
of the unipotents.  It proves class inversion on the cover and on
its central quotient; compare the explicit class action in
\cite[Lemma~3.1, p.~7]{BI:WhiteSL2}.

If the irreducible character $\chi$ of $E$ is real, D-BI-18a
again gives a contradiction.  If it is nonreal, the nontrivial
diagonal outer automorphism sends $\chi$ to
$\overline\chi\ne\chi$.  It cannot be induced by conjugation
by an element of $J$, since such conjugation preserves the
character of the restriction of the fixed representation on $V$.
Hence every element of $J$ induces an inner automorphism of $E$.
For $g\in J$, choose $x\in E$ inducing the same automorphism.
Then $gx^{-1}$ centralizes the irreducible group $E$, and is
scalar by Schur's lemma.  Consequently
\[
                             J\subset\mathbf C^*E.
\]
The scalar-only part of D-BI-18 excludes this case because
$q=31\ge29$.  This finishes all six fields in $\mathcal Q$ and
proves (BI.72).
\end{proof}

\paragraph{Certificates and scope.}
The file \texttt{check\_lie\_small\_fields.py} verifies the exact
ceiling sets, the total-degree comparisons, the stated exceptional
minimum-degree expressions, and the $q=37$ disjointness.  It does
not reprove the published degree tables, classify all simple groups,
or test period dominance by a dimension inequality.  The
inverse-conjugacy and automorphism arguments above are universal
two-by-two proofs, not conclusions drawn from finite samples.
The remaining Lie fields are $q\le27$; alternating cores and
other simple core types are outside this theorem's scope.
